% xelatex -no-pdf pam-1.tex
% dvisvgm --no-fonts --exact --bbox=papersize --output=pam-1.svg pam-1.xdv
\def\pgfsysdriver{pgfsys-dvisvgm.def}
\documentclass[tikz,border=5pt]{standalone}
\usepackage[UTF8,fontset=fandol]{ctex}
\usetikzlibrary{arrows.meta}
\begin{document}
\begin{tikzpicture}[
    line width=0.8pt,
    state/.style={circle,draw,fill=cyan!45,minimum size=1.6cm,inner sep=1pt,font=\Large},
    suffix/.style={dashed,-{Stealth[length=2.5mm]},shorten >=1pt}
  ]
  \node[font=\large] at (2.65,1.3) {字符串 \texttt{abba} 构建的回文树};
  \node[state] (even) at (0,0) {$0$};
  \node[state] (odd) at (4,0) {$-1$};
  \node[state] (bb) at (0,-2.5) {\texttt{bb}};
  \node[state] (a) at (3,-2.5) {\texttt{a}};
  \node[state] (b) at (5.3,-2.5) {\texttt{b}};
  \node[state] (abba) at (0,-5) {\texttt{abba}};
  % Solid edges add the same character at both ends of a palindrome.
  \draw (even) -- (bb) -- (abba);
  \draw (odd) -- (a);
  \draw (odd) -- (b);
  % Each dashed edge points to the longest proper palindromic suffix.
  \draw[suffix] (even) -- (odd);
  \draw[suffix] (a.north west) to[out=145,in=-35] (even.south east);
  \draw[suffix] (b.north west) to[out=145,in=-10] (even.south east);
  \draw[suffix] (bb.south east) to[out=-30,in=-135] (b.south west);
  \draw[suffix] (abba.north east) to[out=30,in=-100] (a.south);
  \node[font=\large] at (4.25,-5) {虚线为后缀链接};
\end{tikzpicture}
\end{document}
